% ══════════════════════════════════════════════════════════════════════════════
%  ADVANCED SCIENCE SKILLS — the above-level book (from 2027)
%  Faculty decision, Nov 2026: every student completes the at-level lesson book;
%  students who excel move on to this book. Its topics go deeper into the ideas
%  of the at-level lessons rather than adding more of the same work.
%
%  Each topic: a reading, the deeper idea worked through (I do / We do), then
%  %%PARTD:N%% pulls in the old Part D questions of at-level Lesson N as
%  practice, then a Challenge. Built by make_a5.py (via a5_levels.py).
% ══════════════════════════════════════════════════════════════════════════════

% ══════════════════════════════════════════════════════════════════════════════
%  TOPIC 1 — HOW SCIENCE CHANGES ITS MIND                    (deepens Lesson 1)
% ══════════════════════════════════════════════════════════════════════════════
\setcounter{section}{0}\renewcommand{\lessonart}{1}
\sectiondiv{How Science Changes Its Mind}{Peer review and replication \textbullet{} Judging a claim \textbullet{} Correlation is not cause}\label{sec:1}
\lessoncodes{Deepens Lesson 1 \textbullet{} VC2S10H01 \textbullet{} VC2S10I06 \textbullet{} VC2S10I07}

\vocabhead
\begin{ctab}
\begin{tabular}{|SW{3.4cm}|W{7.0cm}|}
  \hline\hd & \thead{Word} & \thead{Definition (copy from the slide)} \\ \hline
  & \scaf{\textbf{Peer review}} & \scaf{Checking of a scientist's work by other experts before it is published.} \\ \hline
  & \textbf{Replicate} & \\ \hline
  & \textbf{Correlation} & \\ \hline
  & \textbf{Cause} & \\ \hline
  & \textbf{Placebo} & \\ \hline
  & \textbf{Blind trial} & \\ \hline
\end{tabular}
\end{ctab}

\lessonpart{A}{Reading and comprehension}{Teal}\label{read:1}

\begin{readingbox}{Reading --- Too Good to Be True}
In 1989 two chemists announced that they had produced nuclear fusion, the energy source of the Sun, in a jar of water on a laboratory bench. ``Cold fusion'' promised endless cheap energy, and it made headlines around the world. Within weeks, laboratories everywhere tried to \textbf{replicate} the experiment: to repeat it and get the same result. Almost none could. Within a year most scientists had concluded the original results were mistaken.

The same checking works in the other direction. When Barry Marshall and Robin Warren in Perth claimed that most stomach ulcers were caused by bacteria, not stress, other scientists doubted them. Marshall even drank a culture of the bacteria and became ill. As other researchers replicated their findings, the idea was accepted, and in 2005 they won the Nobel Prize.

Before a study is published, other experts check its methods and reasoning: this is \textbf{peer review}. It cannot catch every mistake, but together with replication it is how science corrects itself. One trap is confusing a \textbf{correlation}, two things that tend to happen together, with a \textbf{cause}, where one thing makes the other happen.
\end{readingbox}

\rcreset
\rcq{What does it mean to replicate an experiment?}{2cm}
\rcq{How did replication settle the cold fusion claim?}{2.5cm}
\rcq{Why was the ulcer idea eventually accepted, even though many doubted it at first?}{2.5cm}
\rcq{What is the difference between a correlation and a cause?}{2.5cm}

\lessonpart{B}{Theory and examples}{Blue}\label{theory:1}

\subsection{Judging a Claim}\label{sub:a1.1}

\begin{skillbox}{Five questions to ask of any claim}
\textbf{1. Who says so,} and do they gain from it? \quad \textbf{2. What is the evidence:} a test, or only stories? \quad \textbf{3. How many} were tested? \quad \textbf{4. Was there a comparison group} that did not get the treatment? \quad \textbf{5. Has anyone else replicated it?}
\end{skillbox}

\begin{workedex}{A memory drink}
A company website says its drink ``improves memory by 30\%''. It gave the drink to 10 of its own staff, who all said they felt sharper.

\smallskip
The company profits if you believe it. The evidence is only how people felt, not a measured test. Ten people is a tiny sample. Nobody drank a plain drink for comparison, and no one else has repeated it. The claim is \textbf{weak} on every question.
\end{workedex}

\begin{wedo}{A study app}
A news report says students who use a study app score higher on tests. The study, by a university, compared 2000 users with 2000 non-users. Use the five questions to judge the claim. What else would you want to know?
\answerbox{5cm}
\end{wedo}

\subsection{Correlation Is Not Cause}\label{sub:a1.2}

\begin{workedex}{Ice cream and shark attacks}
On days when more ice cream is sold, there are more shark attacks. Does ice cream cause shark attacks?

\smallskip
No. Both go up on hot days, when more people buy ice cream \emph{and} more people swim. A third factor, the weather, causes both. The two are correlated, but one does not cause the other.
\end{workedex}

\begin{wedo}{Bigger shoes, better spelling}
Across a primary school, students with bigger shoe sizes are better at spelling. Explain why, and name the third factor.
\answerbox{4cm}
\end{wedo}

\lessonpart{C}{Practice}{Amber}\label{prac:1}

%%PARTD:1%%

\lessonpart{D}{Challenge}{Violet}\label{ext:1}

\task A company has a new cream it says eases sore muscles. Design a fair trial. Explain how you would use a \textbf{placebo} (a cream that looks and smells the same but has no active ingredient), why the people in the trial should not know which cream they have (a \textbf{blind trial}), and why it is even better if the researchers checking them do not know either.
\answerbox{9cm}

\clearpage

% ══════════════════════════════════════════════════════════════════════════════
%  TOPIC 2 — DESIGNING BETTER EXPERIMENTS                    (deepens Lesson 4)
% ══════════════════════════════════════════════════════════════════════════════
\setcounter{section}{1}\renewcommand{\lessonart}{4}
\sectiondiv{Designing Better Experiments}{Repeats and averages \textbullet{} Outliers \textbullet{} Accuracy and precision \textbullet{} Making results trustworthy}\label{sec:2}
\lessoncodes{Deepens Lesson 4 \textbullet{} VC2S10I02 \textbullet{} VC2S10I03 \textbullet{} VC2S10I05}

\vocabhead
\begin{ctab}
\begin{tabular}{|SW{3.4cm}|W{7.0cm}|}
  \hline\hd & \thead{Word} & \thead{Definition (copy from the slide)} \\ \hline
  & \scaf{\textbf{Trial}} & \scaf{One run of an experiment.} \\ \hline
  & \textbf{Mean} & \\ \hline
  & \textbf{Outlier} & \\ \hline
  & \textbf{Accurate} & \\ \hline
  & \textbf{Precise} & \\ \hline
  & \textbf{Reliable} & \\ \hline
\end{tabular}
\end{ctab}

\lessonpart{A}{Reading and comprehension}{Teal}\label{read:2}

\begin{readingbox}{Reading --- Once Is Never Enough}
Drop a paper helicopter and time its fall: 2.31 seconds. Drop it again: 2.28. Again: 2.35. No two \textbf{trials} give quite the same answer, because no two drops are quite the same: your reaction time, a draught, the angle of release. That is why scientists repeat measurements and take the \textbf{mean}, the average. Random errors in one direction tend to cancel errors in the other.

Sometimes one result is far from the rest: 3.10 seconds, perhaps because the stopwatch was started late. A result that does not fit the pattern is an \textbf{outlier}. Scientists look for a reason, and if it was a mistake, leave it out of the mean and say so.

Two other words matter. A measurement is \textbf{accurate} if it is close to the true value. Measurements are \textbf{precise} if they are close to each other. Kitchen scales that always read 50\,g too heavy are precise but not accurate. Results that come out much the same when the experiment is repeated, by you or someone else, are \textbf{reliable}.
\end{readingbox}

\rcreset
\rcq{Why do repeated trials give slightly different results?}{2.5cm}
\rcq{Why does taking a mean give a better answer than one trial?}{2.5cm}
\rcq{What should a scientist do about an outlier?}{2.5cm}
\rcq{Explain how scales can be precise but not accurate.}{2.5cm}

\lessonpart{B}{Theory and examples}{Blue}\label{theory:2}

\subsection{Means and Outliers}\label{sub:a2.1}

\begin{workedex}{Helicopter times}
Times (s): 2.31, 2.28, 2.35, 3.10, 2.30.

\smallskip
3.10 is far from the others, which are all within 0.07\,s of each other, so it is an \textbf{outlier} (probably a late start). Leave it out:\\
mean $= (2.31 + 2.28 + 2.35 + 2.30) \div 4 = 9.24 \div 4 = \mathbf{2.31\ s}$.
\end{workedex}

\begin{wedo}{Ball bounce heights}
Bounce heights (cm): 42, 44, 29, 43. Identify the outlier, suggest a reason for it, and calculate the mean without it.
\workbox{4.5cm}
\end{wedo}

\subsection{Accurate and Precise}\label{sub:a2.2}

\begin{workedex}{Two thermometers}
Boiling water (100\,\textdegree C) is measured three times with each of two thermometers.
A: 100.1, 99.9, 100.0. \quad B: 96.0, 96.1, 95.9.

\smallskip
A is accurate and precise. B is precise (its readings agree) but not accurate (they are all about 4\,\textdegree C low). B has a \textbf{systematic error}: the same error every time, which repeating cannot fix, but checking against a known value can.
\end{workedex}

\begin{wedo}{Two more thermometers}
C: 98.0, 102.5, 99.5. \quad D: 91.0, 97.0, 94.5. Describe each as accurate or not, and precise or not. Which would you trust least, and why?
\answerbox{4.5cm}
\end{wedo}

\lessonpart{C}{Practice}{Amber}\label{prac:2}

%%PARTD:4%%

\lessonpart{D}{Challenge}{Violet}\label{ext:2}

\task A student tests whether fertiliser makes bean plants grow taller. She puts one plant with fertiliser on a sunny windowsill and one without on a shady shelf, waters them ``when they look dry'', and measures each once, after a week, with a ruler that starts 5\,mm from the end. She concludes fertiliser makes plants grow taller. Identify four problems with her method and explain how you would fix each.
\answerbox{9cm}

\clearpage

% ══════════════════════════════════════════════════════════════════════════════
%  TOPIC 3 — MEASURING AND CONTROLLING RISK                  (deepens Lesson 5)
% ══════════════════════════════════════════════════════════════════════════════
\setcounter{section}{2}\renewcommand{\lessonart}{5}
\sectiondiv{Measuring and Controlling Risk}{Hazard and risk \textbullet{} The risk matrix \textbullet{} The hierarchy of controls \textbullet{} Writing a risk assessment}\label{sec:3}
\lessoncodes{Deepens Lesson 5 \textbullet{} VC2S10I02 \textbullet{} VC2S10I03}

\vocabhead
\begin{ctab}
\begin{tabular}{|SW{3.4cm}|W{7.0cm}|}
  \hline\hd & \thead{Word} & \thead{Definition (copy from the slide)} \\ \hline
  & \scaf{\textbf{Likelihood}} & \scaf{How probable it is that harm will happen.} \\ \hline
  & \textbf{Consequence} & \\ \hline
  & \textbf{Risk matrix} & \\ \hline
  & \textbf{Hierarchy of controls} & \\ \hline
  & \textbf{Residual risk} & \\ \hline
\end{tabular}
\end{ctab}

\lessonpart{A}{Reading and comprehension}{Teal}\label{read:3}

\begin{readingbox}{Reading --- How Safety Officers Think}
A hazard is anything that could cause harm. Risk is the chance that it actually will, and how bad it would be. Safety officers in laboratories, factories and hospitals judge risk with two questions. What is the \textbf{likelihood} that someone is harmed? And what is the \textbf{consequence}, how serious would the harm be? A \textbf{risk matrix} combines the two answers into a single rating.

Not all ways of reducing risk are equally good. The \textbf{hierarchy of controls} ranks them. Best is to remove the hazard altogether. Next, swap it for something less dangerous. Then use equipment that separates people from the hazard, such as a guard on a machine or a fume cupboard. Below that come rules, signs and training, which only work if people follow them. Last, and least effective, is protective clothing, such as gloves and safety glasses, which only protects the person wearing it, and only if it is worn properly.

After the controls are in place, some risk always remains. This \textbf{residual risk} must be low enough to accept before the work goes ahead.
\end{readingbox}

\rcreset
\rcq{What two questions does a safety officer ask to judge a risk?}{2.5cm}
\rcq{Why is removing a hazard better than wearing gloves?}{2.5cm}
\rcq{Why are rules and signs lower on the hierarchy than a guard on a machine?}{2.5cm}
\rcq{What is residual risk?}{2cm}

\lessonpart{B}{Theory and examples}{Blue}\label{theory:3}

\subsection{The Risk Matrix}\label{sub:a3.1}

\begin{ctab}
\begin{tabular}{|C{2.6cm}|C{2.4cm}|C{2.4cm}|C{2.4cm}|}
  \hline & \thead{Minor (1)} & \thead{Moderate (2)} & \thead{Serious (3)} \\ \hline
  \thead{Likely (3)} & 3 medium & 6 high & 9 high \\ \hline
  \thead{Possible (2)} & 2 low & 4 medium & 6 high \\ \hline
  \thead{Unlikely (1)} & 1 low & 2 low & 3 medium \\ \hline
\end{tabular}
\end{ctab}
\begin{center}\small Risk rating $=$ likelihood $\times$ consequence\end{center}

\begin{workedex}{Heating water with a Bunsen burner}
Hazard: a scald from boiling water. Without controls a spill is \textbf{possible} (2) and a scald is \textbf{moderate} (2): $2 \times 2 = 4$, medium.
With controls (beaker on a tripod and gauze, standing up, hot water left to cool before it is moved), a spill becomes \textbf{unlikely} (1): $1 \times 2 = 2$, low. The residual risk is acceptable.
\end{workedex}

\begin{wedo}{Cutting with a scalpel}
Hazard: cutting a finger while dissecting a flower. Rate the risk before and after controls, and name the controls you would use.
\answerbox{5cm}
\end{wedo}

\subsection{The Hierarchy of Controls}\label{sub:a3.2}

\begin{skillbox}{Most to least effective}
\textbf{1.} Eliminate the hazard \quad \textbf{2.} Substitute something safer \quad \textbf{3.} Engineering controls (guards, fume cupboards) \quad \textbf{4.} Administrative controls (rules, training, signs) \quad \textbf{5.} Personal protective equipment (gloves, glasses, coats)
\end{skillbox}

\begin{workedex}{A strong-smelling solvent}
An experiment uses a solvent whose fumes cause headaches. \textbf{Substitute} a safer solvent if one works; if not, use it in a \textbf{fume cupboard} (engineering), \textbf{train} students and limit the amount (administrative), and wear \textbf{gloves and safety glasses} (PPE). Several controls together give the lowest residual risk.
\end{workedex}

\begin{wedo}{Broken glass}
Students keep cutting themselves on chipped beakers. Suggest one control from each of three different levels of the hierarchy, and say which you think would work best.
\answerbox{5cm}
\end{wedo}

\lessonpart{C}{Practice}{Amber}\label{prac:3}

%%PARTD:5%%

\lessonpart{D}{Challenge}{Violet}\label{ext:3}

\task Write a full risk assessment for heating salt water to dryness in an evaporating basin over a Bunsen burner. For three hazards, give the likelihood and consequence before controls, the controls (with their level in the hierarchy), and the risk rating after controls.
\begin{ctab}
\begin{tabular}{|W{2.0cm}|C{1.3cm}|W{3.6cm}|C{1.3cm}|}
  \hline\thead{Hazard} & \thead{Risk before} & \thead{Controls (and level)} & \thead{Risk after} \\ \hline
  & & & \\[12mm] \hline
  & & & \\[12mm] \hline
  & & & \\[12mm] \hline
\end{tabular}
\end{ctab}

\clearpage
